\section{Related Work}

\textbf{Multi-Agent LLM Systems.}
Recent work has explored decomposing complex tasks into specialized agents. LangGraph~\cite{langgraph} provides a framework for orchestrating multi-agent workflows using state machines, while systems like AutoGPT~\cite{autogpt} and MetaGPT~\cite{metagpt} demonstrate emergent capabilities through agent collaboration. These approaches share the principle of task decomposition but differ in coordination mechanisms and specialization strategies. Our work extends this paradigm to explanation generation, introducing staged batching for computational efficiency and adaptive synthesis for quality-aware content mixing. Unlike prior systems that primarily focus on tool use or code generation, we address the unique challenge of balancing accuracy and simplicity in educational content.

\textbf{Retrieval-Augmented Generation.}
RAG approaches~\cite{lewis2020rag} enhance language model capabilities by grounding generation in retrieved documents. Hybrid retrieval strategies combining sparse (BM25) and dense (neural) methods have shown improvements over single-strategy approaches~\cite{ma2021hybrid}. Question-answering systems~\cite{izacard2021fid} leverage retrieval to reduce hallucinations and improve factual accuracy. Our multi-agent architecture integrates hybrid RAG (BM25 + semantic search over OpenThoughts-114k~\cite{openthoughts} and Wikipedia) within a structured pipeline, enabling explicit fact extraction and quality-based synthesis rather than end-to-end generation.

\textbf{ELI5 and Evaluation Frameworks.}
The ELI5 dataset~\cite{fan2019eli5} provides a benchmark for simplified explanations, consisting of complex questions paired with community-generated answers. Evaluation of generated explanations remains challenging due to the open-ended nature of the task. The RAGAS framework~\cite{ragas} provides LLM-based metrics for assessing RAG systems, including answer accuracy and faithfulness. Traditional metrics like ROUGE~\cite{lin2004rouge}, BLEU~\cite{papineni2002bleu}, and BERT-Score~\cite{zhang2020bertscore} capture lexical and semantic similarity. We employ both LLM-based and automatic metrics to provide comprehensive evaluation, revealing that these metric families can yield conflicting assessments of explanation quality---a finding with important implications for evaluation methodology in this domain.
